\begin{pastexam}[label={px:ccu-109-swa-01}]{中正資工所甲組 109 學年度 第 4 題}
（程式除錯）

(20\%) Below is a \textbf{bubble sort} program that sorts an integer array containing 10 elements.
The algorithm sorts elements in \textbf{ascending order}, i.e., the greatest element in the list
will be carried to the end of the list. However, the first version of the program contains a lot
of bugs. Please indicate and correct all bugs so that the program creates the desired output.

Try to use as \textbf{few changes} as possible to make the program compile and run correctly. You
may also insert and delete lines if you like to. Do not rewrite entire lines of code but try to
keep the changes as \textbf{small} as possible.
\begin{minted}[fontsize=\footnotesize]{text}
Line 1:   #include <iostream>
Line 2:   #include <vector>
Line 3:   using namespace std;
Line 4:
Line 5:   void SwapIntegers (int a, int b) {
Line 6:       int temp = a;
Line 7:       a = b;
Line 8:       temp = b; }
Line 9:
Line 10:  void BubbleSort (vector &intVector) {
Line 11:      for (int i = intVector.size() - 1; i > 0; i++)
Line 12:          for (j = 0; j < i; j++)
Line 13:              if (intVector[j] < intVector[j + 1])
Line 14:                  SwapIntegers (intVector[j], intVector[j + 1]); }
Line 15:
Line 16:  int main () {
Line 17:      int intArray[] = {34, 1432, 1, -54, 16, 22, 13245, 512, -3000, 0};
Line 18:      vector<int> intVector(intArray, intArray + 9);
Line 19:
Line 20:      BubbleSort (&intVector);
Line 21:
Line 22:      for (int i = 0; i <= intVector.size(); i++)
Line 23:          cout >> intVector[i] >> endl;
Line 24:
Line 25:      return 0; }
\end{minted}
\end{pastexam}

\begin{pastsoln}
共 $10$ 處要改，都只動一小段。 先看程式想做什麼：外層 $i$ 從最後一格往前，每一輪內層 $j$ 從
頭掃到 $i - 1$，相鄰兩個前大後小就交換，一輪結束時這一段最大的數被推到第 $i$ 格。 這就是資結書
虛擬碼~14.1 的氣泡排序（資結書第~14.2~節），只是索引從 $0$ 開始。 照這個意思逐行檢查：

\begin{center}
\small
\begin{tabular}{rp{0.48\linewidth}l}
\toprule
行 & 錯在哪裡 & 改的地方 \\
\midrule
5 & 傳值呼叫，函式裡交換的是副本，呼叫端的元素不變 & \texttt{int \&a, int \&b} \\
8 & 沒有把 \texttt{temp} 存的舊 \texttt{a} 放回 \texttt{b} & \texttt{b = temp;} \\
10 & \texttt{vector} 是樣板，要寫元素型別，編譯錯誤 & \texttt{vector<int> \&intVector} \\
11 & $i$ 應該往下數；\texttt{i++} 讓 $i$ 越來越大，迴圈不會停 & \texttt{i-{}-} \\
12 & \texttt{j} 沒有宣告，編譯錯誤 & \texttt{int j = 0} \\
13 & 前小後大就交換，會把最小的推到最後，排成遞減 & \texttt{intVector[j] > intVector[j + 1]} \\
18 & 只放了前 $9$ 個數，漏掉最後的 $0$ & \texttt{intArray + 10} \\
20 & 參數是參考，直接傳物件；\texttt{\&intVector} 是指標，型別不合，編譯錯誤 & \texttt{BubbleSort (intVector)} \\
22 & 用 \texttt{<=} 會讀到 \texttt{intVector[10]}，超出範圍 & \texttt{i < intVector.size()} \\
23 & 輸出要用 \texttt{<<}，\texttt{>>} 是輸入運算子，編譯錯誤 & \texttt{cout << intVector[i] << endl} \\
\bottomrule
\end{tabular}
\end{center}

幾點說明：
\begin{itemize}
\item 第 18 行的建構子 \texttt{vector<int>(first, last)} 複製 $[\mathit{first}, \mathit{last})$
這段，不含 \texttt{last}。 \texttt{intArray + 9} 只拿到第 $0$ 到第 $8$ 個數，題目說要排 $10$ 個，
所以改成 \texttt{+ 10}。
\item 第 13 行：題目要求遞增，而且「最大的被帶到最後」。 原來的 \texttt{<} 在前面比較小時交換，
每一輪把最小的帶到最後，結果是遞減，所以改成 \texttt{>}。
\item 第 5、8 行兩個錯要一起改：只加 \texttt{\&} 而不改第 8 行，\texttt{a}、\texttt{b} 都會變成原來的
\texttt{b}；只改第 8 行而不加 \texttt{\&}，呼叫端完全不變。
\item 第 11 行 \texttt{intVector.size() - 1} 是無號數 $9$，指定給 \texttt{int} 沒有問題（陣列不是空的）。
第 22 行改完以後 \texttt{int} 和無號數比較，編譯器可能給警告，但不影響正確性，不必改。
\item 第 14 行最後的 \texttt{\}} 是 \texttt{BubbleSort} 的函式結尾。 兩層 \texttt{for} 都只有一個
敘述，不加大括號也對。
\end{itemize}

改完的整支程式（左邊的行號就是原題的 Line 編號）：
\begin{minted}[fontsize=\footnotesize, linenos, numbersep=6pt, style=perldoc]{cpp}
#include <iostream>
#include <vector>
using namespace std;

void SwapIntegers (int &a, int &b) {
    int temp = a;
    a = b;
    b = temp; }

void BubbleSort (vector<int> &intVector) {
    for (int i = intVector.size() - 1; i > 0; i--)
        for (int j = 0; j < i; j++)
            if (intVector[j] > intVector[j + 1])
                SwapIntegers (intVector[j], intVector[j + 1]); }

int main () {
    int intArray[] = {34, 1432, 1, -54, 16, 22, 13245, 512, -3000, 0};
    vector<int> intVector(intArray, intArray + 10);

    BubbleSort (intVector);

    for (int i = 0; i < intVector.size(); i++)
        cout << intVector[i] << endl;

    return 0; }
\end{minted}
執行結果是一行一個數：$-3000, -54, 0, 1, 16, 22, 34, 512, 1432, 13245$。
\end{pastsoln}
